% ОБЯЗАТЕЛЬНО ИМЕННО ТАКОЙ documentclass!
% (Основной кегль = 14pt, поэтому необходим extsizes)
% Формат, разумеется, А4
% article потому что стандарт не подразумевает разделов
% Глава = section, Параграф = subsection
% (понятия "глава" и "параграф" из документа, описывающего диплом)
\documentclass[a4paper,article,14pt]{extarticle}

% Подключаем главный пакет со всем необходимым
\usepackage{thesis}

% Пакеты по желанию (самые распространенные)
% Хитрые мат. символы
\usepackage{euscript}
% Таблицы
\usepackage{longtable}
\usepackage{makecell}
% Картинки (можно встявлять даже pdf)
\usepackage[pdftex]{graphicx}

\usepackage{amsthm,amssymb, amsmath}
\usepackage{textcomp}


\begin{document}

% Титульный лист диплома СПбГУ
% Временное удаление foot на titlepage
\newgeometry{left=30mm, top=20mm, right=15mm, bottom=20mm, nohead, nofoot}
\begin{titlepage}
\begin{center}

\textbf{Санкт--Петербургский}
\textbf{государственный университет}

\vspace{35mm}

\textbf{\textit{\large ТРЕТЬЯКОВ Даниил Александрович}} \\[8mm]
% Название
\textbf{\large Выпускная квалификационная работа}\\[3mm]
\textbf{\textit{\large Исследование эффективности непрерывного обучения с использованием сетей Колмогорова-Арнольда в задаче распознавания образов}}

\vspace{20mm}
Уровень образования: бакалавриат\\
Направление 01.03.02 «Прикладная математика и информатика»\\
Основная образовательная программа СВ.5005.2021
«Прикладная математика, фундаментальная информатика и программирование»\\[25mm]


% Научный руководитель, рецензент
\begin{flushright}
\begin{minipage}[t]{0.65\textwidth}
{Научный руководитель:} \\
профессор кафедры компьютерного моделирования и многопроцессорных систем, д.т.н. Дегтярев Александр Борисович

\vspace{10mm}

{Рецензент:} \\
TODO: обновить инфу
\end{minipage}
\end{flushright}

\vfill 

{Санкт-Петербург}
\par{2025 г.}
\end{center}
\end{titlepage}
% Возвращаем настройки geometry обратно (то, что объявлено в преамбуле)
\restoregeometry
% Добавляем 1 к счетчику страниц ПОСЛЕ titlepage, чтобы исключить 
% влияние titlepage environment
\addtocounter{page}{1}

% Содержание
\tableofcontents
\pagebreak

\specialsection{Введение}

В последние десятилетия нейронные сети стали неотъемлемой частью современных 
интеллектуальных систем. Их способность аппроксимировать сложные зависимости и 
выявлять скрытые паттерны в больших объемах данных делает их незаменимыми 
инструментами в сложных исследовательских и прикладных задачах. Именно они
сделали возможным появление алгоритмов, превосходящих человеческие возможности в
областях, где люди традиционно считались непревзойденными, например — в 
распознавании изображений и речи, играх и стратегиях, медицинской диагностике.

Однако после того, как нейронная сеть была обучена выполнению определенной 
задачи, например, классификации деревьев, её обучение дополнительной задаче, 
например, классификации птиц, оказывается затруднительным. Это связано с тем, 
что при добавлении новой задачи типичная нейронная сеть склонна\, "забывать"\, 
предыдущую. В связи с этим, одним из актуальных направлений в развитии нейронных
сетей является непрерывное обучение - подход, при котором модель обучается и 
адаптируется на протяжении всего времени её эксплуатации с целью поддержания 
своей производительности в условиях изменяющихся данных и окружающей среды, без 
необходимости полного переобучения. Потребность в подобных моделях особенно 
сильна в рекламной индустрии, финансовой аналитике, робототехнике и 
кибербезопасности.

Поскольку феномен катастрофического забывания является основным вызовом в 
области непрерывного обучения, для борьбы с ним разработаны подходы, которые 
можно разделить на четыре категории:
\begin{enumerate}
    \item С использованием ансамблирования: предполагается построение для каждой 
    новой задачи специализирующейся на ней модели с последующим объединением 
    всех моделей в единый ансамбль.
    \item С использованием регуляризации: в функцию потерь добавляются 
    регуляризирующие члены, призванные сохранить важные параметры, связанные с 
    предыдущими задачами.
    
    \pagebreak % Prevents the 3rd item from being devided into two pages
    
    \item С использованием разреженного кодирования представлений: разреженность
    в представлениях подразумевает, что только небольшая часть нейронов 
    активируется для обработки конкретного входного сигнала, что снижает 
    вероятность "перезаписи" информации.
    \item С воспроизведением опыта: основаны на идее периодического повторения 
    данных из предыдущих задач во время обучения на новых задачах. Это позволяет 
    модели "освежать" свои знания о старых задачах и предотвращать их забывание.
\end{enumerate}

Как правило, в исследованиях, посвященных методам борьбы с катастрофическим 
забыванием, чаще всего рассматривается одна и та же архитектура - многослойный 
перцептрон (Multi-Layer Perceptron, MLP). Это связано с тем, что MLP - одна из 
самых простых и универсальных архитектур нейронных сетей - обладает 
основополагающим свойством: способностью выступать в роли универсального 
аппроксиматора, что позволяет эффективно обрабатывать высокоуровневые признаки. 
Благодаря этому MLP становится важным компонентом многих современных моделей, 
таких как CNN, RNN и трансформеры.

Несмотря на определённые успехи в борьбе с катастрофическим забыванием в 
многослойных перцептронах, существует альтернативная архитектура, подходящая для 
использования в качестве универсального аппроксиматора — сети 
Колмогорова-Арнольда (Kolmogorov-Arnold Networks, KAN). В отличие от MLP, где 
функции активации фиксированы и применяются к выходам нейронов, в KAN функции 
активации применяются на месте традиционных линейных весов и являются 
обучаемыми. Они могут быть параметризованы, например, с помощью базисных 
сплайнов. Такая архитектура основывается на теореме Колмогорова-Арнольда, 
которая утверждает, что любая непрерывная функция может быть представлена в виде 
суперпозиции унарных функций и бинарной операции сложения.

Теоретическая основа KAN, в частности, свойство локальности базисных сплайнов, 
позволяет более эффективно контролировать влияние отдельных параметров на 
итоговое представление, что, в свою очередь, может способствовать снижению 
проявлений катастрофического забывания. Поэтому ожидается, что подходы, 
разработанные для борьбы с этим феноменом в MLP, продемонстрируют лучшие 
результаты при их применении в KAN.

Данная работа посвящена разработке и адаптации методов преодоления 
катастрофического забывания для архитектуры сетей Колмогорова-Арнольда, а также 
исследованию эффективности полученных моделей в задачах распознавания образов. В 
рамках исследования планируется провести сравнительный анализ 
усовершенствованных моделей KAN и аналогичных моделей на базе MLP, что позволит 
лучше оценить потенциал новой архитектуры в контексте непрерывного обучения.

\pagebreak

\specialsection{Постановка задачи}

Рассматривается процесс непрерывного обучения нейронных сетей на базе архитектур KAN и MLP в задаче распознавания образов. Датасет \(D\) состоит из размеченных данных, разбитых на \(T\) частей, соответствующих обучающим сессиям, то есть \(D = \{B_t\}^T_{t=1}\).

\subsection*{Цель и задачи исследования}

Целью настоящего исследования является адаптация методов непрерывного обучения в задаче распознавания образов, разработанных для архитектуры MLP, к работе с сетями Колмогорова-Арнольда и последующий анализ полученных с помощью применения адаптированных методов новых архитектур как между собой (в рамках одной архитектуры KAN), так и с соответствующими аналогами на базе архитектуры MLP.

Для достижения поставленной цели необходимо решить следующие задачи:

\begin{itemize}
    \item Провести обзор и анализ существующих методов преодоления катастрофического забывания, применяемых в многослойных перцептронах (MLP).
    
    \item Адаптировать выбранные методы непрерывного обучения для архитектуры сетей Колмогорова-Арнольда (KAN).
    
    \item Разработать усовершенствованные модели KAN с применением адаптированных методов непрерывного обучения.
    
    \item Организовать и провести эксперименты по распознаванию образов с использованием разработанных моделей КАН и аналогичных моделей MLP.
    
    \item Сравнить результаты работы усовершенствованных моделей КАН и MLP для определения их эффективности в условиях непрерывного обучения и проявления катастрофического забывания.
    
    \item Проанализировать полученные экспериментальные данные и сделать выводы о потенциале использования сетей Колмогорова-Арнольда в задачах непрерывного обучения.
\end{itemize}

\pagebreak

\specialsection{Обзор литературы}

\pagebreak

\section{Анализ существующих подходов к преодолению катастрофического забывания}

Здесь будут рассмотрены четыре основных подхода к изменению архитектуры MLP, используемых для смягчения влияния катастрофического забывания. Для каждого из подходов будут описаны представители, сравнение эффективности которых с аналогами на базе архитектуры KAN будет произведено в главе "анализ экспериментов".

\subsection{Подход с использованием ансамблирования}

\subsection{Подход с использованием регуляризации}

\subsection{Подход с использованием разреженного кодирования представлений}

\subsection{Подход с воспроизведением опыта}

\pagebreak

\section{Построение усовершенствованных моделей на базе архитектуры KAN}
Здесь будут описаны изменения, внесённые в архитектуру KAN в соответствии с идеями каждого из рассмотренных в предыдущей главе подходов, а также описаны параметры и гиперпараметры, использованные в каждой из полученных моделей, так, чтобы стало возможным полное вопроизведение процесса усовершенствования.

\subsection{Использование оригинальной архитектуры KAN}

\subsection{Подход с использованием ансамблирования}

\subsection{Подход с использованием регуляризации}

\subsection{Подход с использованием разреженного кодирования представлений}

\subsection{Подход с произведением опыта}

\pagebreak

\section{Применение усовершенствованных моделей в задачах непрерывного обучения}

Здесь будут описаны детали эксперимента с перестановкой данных ( используется для оценки способности модели справляться с катастрофическим забыванием в условиях, когда данные изменяются, но сохраняют ту же структуру) и обучения с постепенным добавлением новых классов, а также приведены таблицы с результатами применения в них каждой из усовершенствованных моделей, полученных в предыдущей главе

\subsection{Обучение с перестановкой данных}

\subsection{Инкрементальное обучение классам}

\pagebreak

\section{Анализ экспериментов}
% \begin{figure}[ht]
% \begin{center}
% \scalebox{0.4}{
%    \includegraphics{images/graph.jpg}
% }

% \caption{
% \label{graph-fig}
%      Линейные функции.}
% \end {center}
% \end {figure}
% Ссылаемся на график ~\ref{graph-fig}.
% Ссылка на статью: \cite{voc}, \cite{vo2}

Здесь будет произведён анализ результатов полученных моделей как между собой (как представителей одной архитектуры), так и с результатами соответствующих аналогов (полученных с использованием тех же подходов) на базе архитектуры MLP.

\subsection{Сравнение усовершенствованных моделей на базе архитектуры KAN}

\subsection{Сравнение эффективности подходов применительно к архитектурам KAN и MLP}

\pagebreak

\specialsection{Выводы}

\pagebreak

\specialsection{Заключение}

\pagebreak

\specialsection{Список литературы}

\pagebreak

\specialsection{Приложение}

% Библиография в cpsconf стиле
% Аргумент {1} ниже включает переопределенный стиль с выравниванием слева
% \begin{thebibliography}{1}
% \bibitem{voc} Griffin D.W., Lim J.S. \flqq Multiband excitation vocoder\frqq. IEEE ASSP-36 (8), 1988, % pp. 1223-1235.
% \bibitem{vo2} Griffin D.W., Lim J.S. \flqq Multiband excitation vocoder\frqq. IEEE ASSP-36 (8), 1988, % pp. 1223-1235.
% \end{thebibliography}

\end{document}